\documentclass[onecolumn, draftclsnofoot,10pt, compsoc]{IEEEtran}
\usepackage{graphicx}
\usepackage{url}
\usepackage{setspace}
\usepackage{listings}

\usepackage{geometry}
\geometry{textheight=9.5in, textwidth=7in}

% 1. Fill in these details
\def \CapstoneTeamName{Manovacumeter Team}
\def \CapstoneTeamNumber{20}
\def \GroupMemberOne{Cade Raichart}
\def \GroupMemberTwo{Daniel Kato}
\def \GroupMemberThree{Nathan Shepherd}
\def \CapstoneProjectName{ISS Barometer App }
\def \CapstoneSponsorCompany{NASA}
\def \CapstoneSponsorPerson{Don Pettit}

% 2. Uncomment the appropriate line below so that the document type works
\def \DocType{Progress Report}

\newcommand{\NameSigPair}[1]{\par
\makebox[2.75in][r]{#1} \hfil 	\makebox[3.25in]{\makebox[2.25in]{\hrulefill} \hfill		\makebox[.75in]{\hrulefill}}
\par\vspace{-12pt} \textit{\tiny\noindent
\makebox[2.75in]{} \hfil		\makebox[3.25in]{\makebox[2.25in][r]{Signature} \hfill	\makebox[.75in][r]{Date}}}}
% 3. If the document is not to be signed, uncomment the RENEWcommand below
%\renewcommand{\NameSigPair}[1]{#1}

%%%%%%%%%%%%%%%%%%%%%%%%%%%%%%%%%%%%%%%
\begin{document}
\begin{titlepage}
    \pagenumbering{gobble}
    \begin{singlespace}
    	\includegraphics[height=4cm]{coe_v_spot1}
        \hfill
        % 4. If you have a logo, use this includegraphics command to put it on the coversheet.
        %\includegraphics[height=4cm]{NASA}
        \par\vspace{.2in}
        \centering
        \scshape{
            \huge CS Capstone \DocType \par
            {\large\today}\par
            \vspace{.5in}
            \textbf{\Huge\CapstoneProjectName}\par
            \vfill
            {\large Prepared for}\par
            \Huge \CapstoneSponsorCompany\par
            \vspace{5pt}
            {\Large\NameSigPair{\CapstoneSponsorPerson}\par}
            {\large Prepared by }\par
            Group\CapstoneTeamNumber\par
            % 5. comment out the line below this one if you do not wish to name your team
            %\CapstoneTeamName\par
            \vspace{5pt}
            {\Large
                \NameSigPair{\GroupMemberOne}\par
                \NameSigPair{\GroupMemberTwo}\par
                \NameSigPair{\GroupMemberThree}\par
            }
            \vspace{20pt}
        }
        \begin{abstract}
        % 6. Fill in your abstract
        The International Space Station, the conglomerate work of multiple national space associations, including NASA, requires precision devices to monitor the many aspects of space life.
        One very important aspect for the success of life in near vacuum is sustained air pressure inside the modules.
        This air pressure is monitored by mechanical barometers, called Manovacometers, which must be visually monitored, and are no longer available.
        The \CapstoneProjectName seeks to provide a suitable replacement for the current barometer device aboard the International Space Station.
        This document will serve to introduce the \CapstoneProjectName, and will detail the problem, the solution, and the success metrics of the project.
        The sections include the Definition of the Problem, Definition of Project, and the Metrics of Success.
        \end{abstract}
    \end{singlespace}
\end{titlepage}
\newpage
\pagenumbering{arabic}
\tableofcontents
% 7. uncomment this (if applicable). Consider adding a page break.
%\listoffigures
%\listoftables
\clearpage

% 8. now you write!
\section{Project Purpose}
The purpose of this project is to create an application, the ISS Barometer, that will be used on the ISS to measure and track air pressure.
The crew members aboard the ISS currently rely on out of date mechanical barometers to take critical readings during leaks.
These leaks, caused by orbital debris, must be found and patched or the crew would need to evacuate.
Our application seeks to improve on the mechanical barometers by providing clear and accurate pressure readings through the iPads on-board barometer sensor.
The app will display the current pressure and changes in pressure, by way of on screen digits and a plot.
Eventually, in version 2, our app will also have the capability of calculating remaining time on vessel before evacuation.
This added functionality will further assist the crew members, allowing them to make quick decisions in potential life and death situations.

\section{Project Goals}
As mentioned above the application we are making will provide ISS crew members with an accurate and efficient pressure reading.
Our required features follow that purpose and are listed below:
\begin{itemize}
\item[V1:] Data Screen with initial pressure, current pressure, start/stop button, time stamp, pressure change (measured in sec / mmHg)
\item[V1:] Graph Screen with plot pressure as a function of time
\item[V1:] Additional settings menu that includes plot scale, scroll functionality, orientation, decimal point, as well as other other units of measure for pressure
\item[V1:] All items will be optimized for efficient reading and clear interface
\item[V2:] Graph will have pinch to zoom functionality, scaling options (to be set in settings)
\item[V2:] Graph will be able to be exported as a csv, excel readable file
\item[V2:] Graph can be compared to t-Res table, to measure time remaining
\item[O1:] Research can be done to identify low pressure accuracy and resistance
\end{itemize}
These features, with the exception of the optional item O1, function as our goals.
Along with these specific features, the application that we create must be efficient and clear, as decided by the users.
The application will be used in challenging situations and must be clear and easy to understand, so that the data can be conveyed quickly and without confusion.

\section{Current State of the Project}
% If anyone can think of a better way to state this, feel free to change the heading
Currently we have an outline of what parts must be included in the app as well as the outline for when the first version of our app must be completed. We have also created our first prototype. This reads from the barometer and displays the current pressure in mm/Hg. It also calculates the dp/dt and displays in below the current pressure.

\section{Project Problems}
We chose to use swift and xcode because they allows for quick development and prototyping.
This said we ran into many problems with that. For one Swift only works on Linux and MacOS.
One of our group members did not have either at the time. They fixed this with installing a linux OS.
Another problem was installing the swift libraries and connecting them with the xcode application.
Linux needed many packages to be downloaded for the library to work correctly, this took many hours of troubleshooting as it was not straightforward.
Connecting swift to xcode on the MacOS was another problem as well.
Our last problem with swift and xcode is we have no experience using either.
We had and still have to continually work on improving our knowledge of these systems using online videos, guides and other resources.

For the first time we all experienced what it was like working with someone who is in the workplace.
We learned that our clients will not be getting back to us immediately and sometimes we have to wait for them reply.
This means we have to work and finish our work and send it to him way earlier then we had planned originally.
We also ran into the problems of finding the correct IEEE documentations.
Sometimes we found the document quick and easy while other times it hours and many students combined to find them.

\section{Interesting Code}
\begin{lstlisting}
  func startDisplayingPressureData() {
      self.altimeter.startRelativeAltitudeUpdates(to: OperationQueue.main, withHandler: {
        data, error in
          self.handlePressureReading(data: data!)
      })
  }

\end{lstlisting}

One of the main features of our application is its ability to read data from the iPad's built in barometer and display the data collected in different ways.
The way that data is collected from the barometer is quite interesting.
The process of reading data from the barometer requires a bit of processing power and time which can block other time sensitive tasks form being performed.
As a result, the measurement is done asynchronously through an object called the \texttt{OperationQueue}.
A method provided by the CoreMotion library called \texttt{startRelativeUpdates()} takes the \texttt{OperationQueue}'s \texttt{main} queue, and a callback function as arguement.
The function then appends a task to the queue to perform in the background, and when the task is complete, the result is passed into the callback function for further processing.\\

\section{Document Summaries}

\subsection{Problem Statement}
The problem statement is where we first expressed the understanding we had of the project.
We described the use case of the barometers that are begin replaced so that our design may fit the needs of the crew members.
This included a definition of the problem which stated that we are designing an iOS application for the iPad Air 2s on the ISS to replace the mechanical barometers called 'MANOVACUMETERS'.
The project was then defined as a three term long project that will result in an application that can be used in place of the MANOVACUMETERS currently being used on the ISS.
This app is to have a graph, a reading of the current pressure in mmHg and dp/dt, and the ability to record pressure by pressing a button.
A settings page is to be included that will allow units, significant digits, plot scale, and orientation.
The success of the project was defined as being when there are two displays; one with initial pressure in mmHg, current pressure in mmHg, pressure change as sec/mmHg and timestamp in GMT, and one with the plot of pressure as a function of time.
The first display will allow the user to start and stop recording of data.

\subsection{Requirements Document}
The Requirements Document is where we discussed details of the project and defined features of the application.
The scope of the project is defined as being to create an application for crew members of the ISS to measure the pressure inside of the ISS in the event of a rapid depressurization event.
Some alternatives to creating a new application were discussed such as using existing barometer apps, but the needs of the space station require a custom application to be built.
The timeline of the project described version 1 being completed by the end of December, and version 2 being completed by mid May.
Each feature is then discussed in detail.
The initial pressure reading feature will show the pressure read when the app loads, or whenever the record button is pressed.
The record button feature will record the pressure readings to be exported as a csv.
The current pressure reading feature will display the current pressure in mmHg.
The time pressure over change feature will display the change in pressure over time as mmHg/sec.
The plot of the pressure change feature will be a graph that displays the change in pressure over time so that the event can be visualized.
The configurable settings page feature will allow the user to change the units, significant digits, orientation, and plot scale.
Then features for version two were described.
The Consult T-Res feature will provide the user with the T-Res measurement that  corresponds to the slope of the graph.
The Save feature will allow the user to save pressure data to a csv.

\subsection{Technology Review}
The goal of this document is to justify the choices made, with research to prove that our project was done in the best way possible.
Included is a list of technology options, details on each, and justifications of each choice.
The nine technologies we needed to find implementations  for was development platform, chart library, UI organization, barometer implementation, data storage, graph library, color palette, programing language and UI library.
We will summaries the appropriate implementations that we chose and the justification for those implementations.

For our development platform technology we will be using Xcode and swift to develop the ISS Barometer application.
We chose this technology for its ability to interface with the onboard barometer, which is required for this project.
The ability to interface directly with the built in barometer is necessary for gaining an accurate pressure reading \cite{techDocKato}.


Our Chart Library Technology will be the library called Charts.
Reliability is of the utmost importance for our application.
Charts offers support for dynamic graphs, while the other 2 libraries don't have the consistent support that will be need.
The Charts API also makes it easy to add points to a dynamic graph, making it the perfect library \cite{techDocKato}.

We are using the single page layout for our Organization of our UI.
Using a single page layout will allow us to access our settings, statistics and our graph frequently and easily.
Since our astronauts will be using this in emergency situations speed and easy access is absolutely necessary \cite{techDocKato}.

Our client has asked us to us the iPad air 2 for our Barometer Implementation technology.
We are asked to use this because the space station has 8-10 iPads aboard.
Using any other implementation would cause us to ship it up to the space station which is not necessary or convenient \cite{techDocRaichart}.

For the Graph Interface technology we have chosen Swift Chars library.
Our client has requested that the graph be able to zoom and pan, the only library that allows us to do this is the swift charts library \cite{techDocRaichart}.

The UI Library will be mostly Cocoa controls for is massive library and larger features list.
Though as we have found no single library has everything we need to complete our task, therefore we will need to use many other libraries such as CRPageViewController and FlatUIKit \cite{techDocRaichart}.


We will be programing in the swift Language.
This was chosen for its ability to quickly prototype and develop the application.
Swift also has many performance benefits for app creation that make it the best option.
Also our client doesn't know much about programming swift works xcode which will make showing him our progress easily \cite{techDocShepherd}.

For our data storage technology we will use CoreData, in tandem with NSUserDefaults.
CoreData allows for large data sets that scale well and far better than the other libraries do.
Using NSUserDefaults package, allows us to quickly save data, such as float point accuracy \cite{techDocShepherd}.

Our color palette technology will be the Subtle Technology Theme.
This was chosen because the other options, OSU similar theme and Nasa similar theme, are designed for specific uses that would be intrusive and counterproductive in our application.
Subtle Technology Theme allows the user to not be distracted by unwanted color \cite{techDocShepherd}.

\subsection{Design Document}
The design document was the final project document, excluding this one.
As the final document, it included all of the research done for the Technology Reviews, as well as previous learnings from our software classes.
Following the IEEE standard given, our team worked together to justify our design choices through various viewpoints.
It was informative to describe the choices we made using these viewpoints, as it required that we make specific arguments to support those choices.
Also it allowed us to incorporate other portions of our project into the detail, justifying not only specific parts, but their interactions as well.
This document was written and edited on the same branch, by all team members.
The viewpoints that we detailed were the context viewpoint, the functional viewpoint, the information viewpoint, the concurrency viewpoint, the development viewpoint, the deployment viewpoint, and post-deployment viewpoint.
With these viewpoints we were able to address the design concerns of all of the stakeholders, and as a result, described the project from many different perspectives.

\section{Retrospective}
\begin{center}
	\begin{tabular}{| l | p{0.3\linewidth} | p{0.3\linewidth} | p{0.3\linewidth} |}
		\hline
		& Positives & Deltas & Actions \\ \hline
    Week 1 & Chose and submitted choices for our projects
      & No changes, just challenges in choosing
      & By justifying our choices, the preferences became clearer \\ \hline
		Week 2 & The group met, we had all chosen the ISS project as our first pick
      & We emailed Don Pettit, for a meeting, and started our Problem Statements
      & We were able to meet on Friday, Don had a compete picture of the app\\ \hline			% Put content between '&'s
		Week 3 & Finished problem statements and turned into both our client and teacher
 & Figure out how to share app with client, whom does not know the github system
 & step by step instructions on how to download the app \\ \hline
		Week 4 & Finalized Problem Statement, got client approval
      & Spelling of MANOVACUMETER was incorrect, among other minor edits
      & Changes were made, and the final Problem Statement was submitted \\ \hline
    Week 5 & Started Requirements Document rough draft & Need to start Requirements Document & Requirements Document will be started \\ \hline
		Week 6 & First prototype with two basic display of barometer numbers
  finished the Requirements Document & & \\ \hline
  	Week 7 & Completed Technology Review document rough draft & Need to complete Technology Review rough draft & A rough draft of the Technology Review will be created \\ \hline
		Week 8 & Started Technology Review & Need to Start Technology Review & Technology Review will be started \\ \hline
		Week 9 &Finished Technology document
Started working on design document
 & Could not find IEEE paper & Searched more and asked around\\ \hline
		Week 10 & Completed Design Document and Progress Report & Need to complete Design Document and Progress Report & Design Document and Progress Report will be completed \\ \hline

	\end{tabular}
\end{center}

\newpage
\bibliographystyle{IEEEtran}
\bibliography{bibliography}
\end{document}
